\section{Discussion}
Observing the analytical and numerical models, angles between \ang{0} to \ang{90} could be observed. Generally, surface wind-turning angles above the classical value of \ang{45} could be observed with strictly decreasing viscosities with altitude. The only case where this did not occur is for a small region in the two-layer model. However, since this region only had a lowest value of \ang{43.4}, the importance of this small region can be dismissed. The reasoning behind decreasing viscosities with altitude providing surface angles over the classical value can be attributed to the decreasing viscosity forcing the rotation due to the Coriolis force to occur at a lower latitude. As the viscosity is initially large but decreases afterwards, the spiral is forced to turn more at the lower altitudes. This has to occur in order to satisfy the geostrophic limit above. Thus, the surface angle becomes larger than the classical value for the decreasing viscosities with altitude.

For the increasing viscosities with altitude, the opposite could be observed with smaller surface wind-turning angles than the classical value. For the two-layer model, only a small region had angles greater than the classical value. However, the largest value was only measured at \ang{46.6}, which is close to the classical limit. The same logic follows to explain why this surface angle became smaller, as the explanation for why the surface angles became larger for the decreasing viscosities with height. If the viscosity increases with height, then the spiral has to turn slower near the surface in order to turn enough to reach the geostrophic limit aloft. And since the viscosity increases with altitude, more turning can occur there, as it is the shear due to the viscosity which forces this turning. 

The same reasoning can be used to explain the limits for the surface angles when the dimensionless layer thickness goes to infinity or zero. Since the dimensionless layer thickness is the ratio between the characteristic height and the Ekman height scale, it follows that when the ratio goes towards infinity, a constant viscosity is obtained, and the angles go to the classical value of \ang{45}. Furthermore, when the ratio between the characteristic height and the Ekman height scale goes towards zero, one would obtain the extremes of the logic presented in the paragraphs above. Since the Coriolis force is perpendicular to the flow, a \ang{90} surface angle will be obtained if the viscosity decreases with altitude, since the entire rotation has to occur at the surface. On the other hand, if the viscosity increases with height, then no part of the rotation can occur at the surface, and a parallel angle of \ang{0} will be obtained at the surface.

Comparing the results from the theoretical models with the ERA5 data, a clear connection can be observed. The ERA5 data showed a clear majority of angles between \ang{0} and \ang{45}, dismissing the background noise in the data. As the atmosphere consists of a turbulent boundary layer, it is reasonable to assume that the viscosity increases with height, which allows for a strong connection between the theoretical models and the ERA5 data. However, there was still a small portion of the data that had angles above \ang{45}, although a boundary layer must be present. This could be attributed to the fact that even strictly increasing viscosities with height can have a small region where the viscosity decreases with height, which would allow for a small region with surface angles above \ang{45}. Although one should not forget that the ERA5 data are reanalysed data, and thus not pure observations, which could explain this portion of the data and the observed noise. However, the bulk of the data is still in agreement with the theoretical models, supporting the findings in this article.

\textbf{COMPARE RESULT IN THIS STUDY WITH THE OTHER STUDIES}
